\subsection{CONTINUITY OF DERIVATIVES}

PROPOSITION 6. Let I be an open interval in R, let \(x_0\) be one of the endpoints of I, and f a vector function defined and continuous on I, with values in a complete normed space E over R; suppose that f has a right derivative at the points of the complement B in I of a countable subset of I. Then for \(\mathbf{f}_d'(x)\) to have a limit as x tends to \(x_0\) while remaining in B and \(\neq x_0\) it is necessary and sufficient that \(\frac{\mathbf{f}(y) - \mathbf{f}(x)}{y - x}\) have a limit c as \((x, y)\) tends to \((x_0, x_0)\) subject to \(x \in I\) , \(y \in I\) , \(x \neq x_0\) , \(y \neq x_0\) and \(x \neq y\) . Under these conditions f extends by continuity to the point \(x_0\) , the right derivative \(\mathbf{f}_{d}^{\prime}(x)\) tends to c as x tends to \(x_{0}\) (while remaining in B) and the function f extended (defined on \(I \cup \{x_{0}\}\) ) has derivative at \(x_{0}\) equal to c.

Suppose for example that \(x_{0}\) is the right-hand endpoint of I. Let us first show that if \(\mathbf{f}_{d}^{\prime}(x)\) tends to c as x tends to \(x_{0}\) while remaining in B and \(\neq x_{0}\) , then \(\frac{\mathbf{f}(y)-\mathbf{f}(x)}{y-x}\) tends to c; this follows immediately from th. 2 applied to the function \(\mathbf{f}(z)-\mathbf{c}z\) , which yields

\[
\| \mathbf {f} (y) - \mathbf {f} (x) - \mathbf {c} (y - x) \| \leqslant (y - x) \sup _ {z \in \mathrm{B}, x <   z <   y} \left\| \mathbf {f} _ {d} ^ {\prime} (z) - \mathbf {c} \right\|
\]

for \(x < y < x_0\) . Conversely, if \(\frac{\mathbf{f}(y) - \mathbf{f}(x)}{y - x}\) tends to \(\mathbf{c}\) , then for every \(\varepsilon > 0\) there exists an \(h > 0\) such that the conditions \(|x - x_0| < h\) , \(|y - x_0| < h\) ( \(x \neq x_0, y \neq x_0\) ) imply

\[
\left\| \mathbf {f} (y) - \mathbf {f} (x) - \mathbf {c} (y - x) \right\| \leqslant \varepsilon | y - x |.\tag{11}
\]

But for all \(x \in B\) and \(\neq x_{0}\) such that \(|x - x_{0}| < h\) there exists a k \textgreater{} 0 (depending on x) such that the relation x \textless{} y \textless{} x + k entails

\[
\left| \left| \mathbf {f} (y) - \mathbf {f} (x) - \mathbf {f} _ {d} ^ {\prime} (x) (y - x) \right| \right| \leqslant \varepsilon | y - x |\tag{12}
\]

from which, considering (11):

\[
\left\| \mathbf {f} _ {d} ^ {\prime} (x) - \mathbf {c} \right\| \leqslant 2 \varepsilon
\]

for \(|x - x_{0}| < h, x \in B\) and \(x \neq x_{0}\) , which proves that \(\mathbf{f}_{d}^{\prime}(x)\) tends to c.~Moreover, from the relation (11) one has immediately that

\[
\left\| \mathbf {f} (y) - \mathbf {f} (x) \right\| \leqslant \left(\left\| \mathbf {c} \right\| + \varepsilon\right) | y - x |,
\]

which proves (by Cauchy's criterion) that \(\mathbf{f}\) has a limit \(\mathbf{d}\) at the point \(x_0\) as \(x\) tends to this point while remaining in \(\mathrm{I}\) and \(\neq x_0\) ; now, letting \(x\) approach \(x_0\) in (11), for \(y \in \mathrm{I}\) , \(y \neq x_0\) and \(|y - x_0| \leqslant h\) , we have

\[
\left\| \frac {\mathbf {f} (y) - \mathbf {d}}{y - x _ {0}} - \mathbf {c} \right\| \leqslant \varepsilon
\]

which proves that \(\mathbf{c}\) is the derivative at the point \(x_0\) of the function \(\mathbf{f}\) extended by continuity to \(\mathrm{I} \cup \{x_0\}\) .

Remark. A similar argument, based on th. 1, shows that if \(f\) is a real function such that \(f_{d}^{\prime}(x)\) tends to \(+\infty\) at the point \(x_0\) then the ratio

\[
(f (y) - f (x)) / (y - x)
\]

also tends to \(+\infty\) , and conversely; if moreover f has a finite limit at the point \(x_{0}\) (which is not a consequence of the present hypothesis), then the function f extended by continuity to \(x_{0}\) has a derivative equal to \(+\infty\) at this point.

